Vietnamese financial phishing arrives mostly through SMS and chat apps such as Zalo, Messenger, Telegram, and Facebook \cite{groupib2022laser,ais2024biometricwarning}. The messages people most want checked are exactly the ones they should not be forwarding anywhere: they contain bank details, phone numbers, one-time passwords, and account-recovery links. This thesis studies a local-first pipeline that fine-tunes a local LLM to detect those threats without sending the message to a cloud service.

\section{Background and Motivation}
The word \textit{Localized} in the title means two things at once, and the project needs both of them.

The first meaning is that the model runs on the user's own device. Analysis happens on a consumer laptop with no cloud API call, so a message full of banking details never leaves the machine. The second meaning is that the model is trained on Vietnamese scam messages specifically, so that it learns Vietnamese banking vocabulary, local scam tactics, and the way real Vietnamese chat messages are actually written. That last point is not decoration. In the training messages, 31.4\% carry chat abbreviations such as \textit{ko}, \textit{tk}, or \textit{acc}, 12.9\% are typed without diacritics at all, so that \textit{tài khoản} appears as \textit{tai khoan}, and 6.3\% contain emoji. A general multilingual model is not trained on any of that.

Neither meaning is enough on its own. A model that runs privately but was never trained on Vietnamese scams would miss the local patterns. A model trained on Vietnamese scams but running in the cloud would still require the user to upload the message it is supposed to protect.

Three requirements shaped the whole project. It handles text only, it runs offline or at least locally by default, and it prioritises catching high-risk messages even at the cost of some false alarms, because missing a real scam is more expensive than raising one extra warning. These requirements make the project smaller than a full fraud assistant that would also handle images and voice calls, but they keep it realistic on a laptop and appropriate for private messages. They also answer a practical safety question: how to reduce dangerous misses without asking the user to trust a score they cannot check.

\section{Problem Statement}
\label{sec:problem-statement}
Two problems motivate this project. Cloud-based fraud checking asks the user to upload a suspicious financial message to a remote service, which creates a privacy risk that does not need to exist. The second problem is that the system has to handle Vietnamese banking vocabulary, local social-engineering tactics, and the way the messages are really written: Vietnamese and English mixed in one sentence, chat abbreviations, missing diacritics, and emoji. Assuming that a general multilingual model already covers all of that is exactly the assumption this project does not make. The target is therefore a local fine-tuning pipeline that classifies a suspicious message on the device and explains its answer, instead of returning a bare confidence score.

The underlying task is supervised multi-class text classification. The input is one pasted message. The output is a structured record with five fields:

\begin{itemize}\itemsep2pt
\item \texttt{label} --- one of four classes: bank impersonation, Zalo social engineering, task scam, or benign;
\item \texttt{risk\_tier} --- benign, suspicious, or high-risk;
\item \texttt{suspicious\_spans} --- the exact pieces of text the decision was based on, quoted from the message itself;
\item \texttt{xai\_explanation} --- a short explanation in Vietnamese of why the message is or is not a scam;
\item \texttt{recommendations} --- safe next steps for the user.
\end{itemize}

\noindent The \texttt{label} field is what the model is graded on, and every training example carries one, assigned when the dataset was built (Section~\ref{sec:quality-check}) rather than added afterwards. The other four fields are produced at the same time, in the same answer. They are extra information for the user; they are not a substitute for the classification.

\section{Research Questions}
Three questions follow from that problem statement:

\begin{enumerate}
\item \textbf{RQ1:} How well do the locally trained Qwen QLoRA and PhoBERT models classify the four Vietnamese threat types, when both are given the same final evaluation?

\item \textbf{RQ2:} Is parameter-efficient fine-tuning and local inference practical on a consumer laptop, and what resource trade-off justifies QLoRA over ordinary LoRA on a GPU with 8~GiB of video memory?

\item \textbf{RQ3:} Can the pipeline return a structured, evidence-based explanation --- risk tier, threat label, suspicious phrases quoted from the message, and safe next steps --- rather than only a confidence score?
\end{enumerate}

\section{Project Contributions}
This project delivers two fine-tuned models together with the dataset they were trained on.

The first contribution is a 2{,}097-message dataset across four classes. Several messages were written from each source article, so two messages built from the same article can differ by little more than a bank name or an amount. All messages from one article are therefore kept together in the same split, so a model is never tested on a message that is a near-copy of one it was trained on.

The second is that the dataset's own faults were found and repaired before any model was trained on it. A review found one whole group of Zalo messages written as third-person narration instead of as real chat messages, and leftover template text spread across the other classes. Those messages were rewritten, dropped, or corrected, and every decision was written down. Chapter~IV describes what was wrong and what was done about it.

The third is training both models locally on exactly the same messages, keeping the full training and validation logs, and exporting the selected Qwen model as a Q8\_0 GGUF file.

The fourth is a single final evaluation of both finished models on the same 220 messages, which also counts how many answers came back unreadable and how many dangerous messages were wrongly called safe.

\section{Scope and Limitations}
The project handles Vietnamese text only. Image-based phishing such as QR codes and forwarded screenshots, voice calls, deepfake audio, and malware analysis are all outside what was built here, although reading text out of a screenshot is proposed as future work in Section~\ref{sec:future-work}.

The dataset is built from public scam-warning articles on \url{https://tinnhiemmang.vn}, and no private user conversations were collected. That choice was made after checking the public alternatives, which are documented in the appendix. Several candidate datasets had no licence, or no clear record of where their data came from, some contained real phone numbers belonging to real people, and one had substantial overlap between its own training and test halves, so any score measured on it would not have been accurate. Public advisory articles avoided all three problems.

The messages in the dataset are generated from those articles rather than collected from real victims, so the dataset is not a substitute for a benchmark built from real traffic.

The dataset was cut into three parts before any training began. The models learn from the \textbf{training} messages. The \textbf{validation} messages are looked at repeatedly during training, to decide which saved version of the model to keep. The third part is \textbf{held out}, which means it was locked away at the moment of the split and neither model, nor any decision in this project, was allowed to see it until everything was finished. It stands in for messages the system has never met, which is the closest a project of this size can get to a real user pasting in a real scam.

Both finished models were then run over the same 220 held-out messages, once. That result was not used to repair data, retrain, move a threshold, or pick a different saved version. Reporting a number from a test set that has been consulted repeatedly is one of the easiest ways to overstate a result, so the single-use rule was fixed before the models were trained rather than defended afterwards.

One honest qualification belongs with that. Automated file-integrity tests had opened the evaluation files before the models were run. Those tests never showed a message to a human, never ran either model, and did not change any result, but it would still be wrong to describe the evaluation set as completely untouched, so this report does not.

The system also analyses one message at a time; following a conversation across several messages is outside the current scope.

\section{Methodology Overview}
The project runs in five stages.

First, scam-warning articles from \textit{tinnhiemmang.vn} were crawled and kept as source material, and Vietnamese phishing messages were generated from them one class at a time. Every message was then checked automatically, scored by a separate model acting as a quality judge, and reviewed by hand where it mattered.

Second, short LoRA and QLoRA test runs measured what the laptop could do. Ordinary LoRA did not crash; it is reported as too demanding on the hardware to be practical.

Third, Qwen3-4B-Instruct-2507 was trained locally with NF4 QLoRA, and PhoBERT was trained locally as a full, non-quantized classifier, both on exactly the same training and validation messages.

Fourth, the chosen Qwen checkpoint was exported as a Q8\_0 GGUF file and load-tested to confirm it works.

Fifth, both finished models were evaluated once on the held-out messages.

Chapter~III explains why each of these steps was designed the way it was, Chapter~IV describes what was actually run, and Chapter~V reports what the runs produced.
